\chapter{Spinori}
Le rappresentazioni di spin esistono per ogni dimensione $D$ dello spaziotempo. Per introdurle, ci si concentra sullo spazio di Minkowski 4-dimensionale $\bb{M}^4$. Il gruppo di Lorentz in quattro dimensioni è $\group{SO}{1,3}$, la cui algebra associata è $\alg{so}{1,3}$. Si può dimostrare che la complessificazione dell'algebra di Lorentz $\alg{so}{1,3}_{\bb{C}} \cong \alg{so}{1,3} \ot \bb{C} $ è isomorfa alla somma diretta di due copie della complessificazione dell'algebra $\alg{su}{2}$, cioè a $\alg{sl}{2,\bb{C}} \cong \alg{su}{2} \ot \bb{C}$, associata al gruppo $\group{SL}{2,\bb{C}}$, ovvero 
\begin{equation}
    \alg{so}{1,3}_{\bb{C}} \cong \alg{sl}{2,\bb{C}} \op \alg{sl}{2,\bb{C}}.
\end{equation}
I generatori di $\alg{su}{2}$ si indicano tipicamente con $T_i$ e soddisfano la relazione di commutazione 
\begin{equation}
    \comm{T_i}{T_j} = if_{ijk}T^k,
\end{equation}
dove le $f_{ijk}$ sono le \emph{costanti di struttura}, che in questo caso coincidono con il simbolo di Levi-Civita $\epsilon_{ijk}$; quindi, più precisamente
\begin{equation}
    \comm{T_i}{T_j} = i\epsilon_{ijk}T^k.
\end{equation}
Una possibile base dell'algebra di Lorentz si può scrivere tramite le matrici di Pauli $\sigma_i$, con le quali si definiscono 
\begin{equation}
    T_i \equiv \frac{\sigma_i}{2},
\end{equation}
dove 
\begin{equation}
    \sigma_1 = 
    \begin{pmatrix}
        0 & 1 \\
        1 & 0
    \end{pmatrix},\quad
    \sigma_2 = 
    \begin{pmatrix}
        0 & -i \\
        i & 0
    \end{pmatrix},\quad
    \sigma_3 = 
    \begin{pmatrix}
        1 & 0 \\
        0 & -1
    \end{pmatrix}.
\end{equation}
Le rappresentazioni irriducibili di dimensione finita del gruppo $\group{SU}{2}$ sono poi etichettate da un numero intero o, eventualmente, semi-intero $j \in \bb{N}_0/2$ con cui si scrive la dimensione della rappresentazione $\rho$, data da
\begin{equation}
    \dim{\rho} = 2j + 1.
\end{equation}
Il fatto interessante è che l'algebra $\alg{sl}{2,\bb{C}}$ possiede relazioni di commutazione della stessa forma di quelle di $\alg{su}{2}$. Si possono pertanto etichettare con un numero $j \in \bb{N}_0/2$ anche le rappresentazioni di $\group{SL}{2,\bb{C}}$. Tuttavia, dato che $\alg{so}{1,3}_{\bb{C}}$ si decompone nella somma diretta di due algebre, sono necessari due numeri $j$ e $j'$ per identificare una rappresentazione del gruppo di Lorentz, le quali vengono appunto indicate con $\bm{(j,j')}$ e la cui dimensione è
\begin{equation}
    \dim{\rho} = (2j + 1)(2j' + 1).
\end{equation}
Le rappresentazioni $\bm{(j,j')}$ con $j + j'$ semi-intero ma non intero sono dette \emph{rappresentazioni di spin}, in quanto ammettono i fermioni\footnote{Più precisamente, le rappresentazioni per cui $j + j'\in \bb{N}_0$ corrispondono a rappresentazioni ordinarie di $\group{SO}{1,3}$, mentre quelle per cui $j + j'$ è semi-intero ma non intero sono rappresentazioni del suo doppio ricoprimento $\group{SL}{2,\bb{C}}$}, a differenza di quelle con $j + j' \in \bb{N}_0$. In particolare, le rappresentazioni di dimensione minore inequivalenti sono due: $\bm{(1/2,0)}$ e $\bm{(0,1/2)}$; entrambe rappresentazioni 2-dimensionali. 

Il prossimo passo consiste nell'omomorfismo tra $\group{SL}{2,\bb{C}}$ e $\group{SO}{1,3}$. Per farlo, ad ogni vettore $x^\mu$ si associa la matrice hermitiana $X$ così definita
\begin{equation}
    x^\mu \mapsto X := 
    \begin{pmatrix}
        x^0 + x^3 & x^1 - ix^2 \\
        x^1 + ix^2 & x^0 - x^3
    \end{pmatrix},
\end{equation}
il cui determinante è $\det{X} = -x_\mu x^\mu$, e si definiscono le matrici
\begin{equation}
\label{eq: matrici di Weyl}
    \sigma_\mu := (-\IM,\sigma_i),\quad \b{\sigma}_\mu \equiv \sigma^\mu := (\IM,\sigma_i),
\end{equation}
dette anche \emph{matrici di Weyl}, che soddisfano le proprietà:
\begin{enumerate}
    \item $\sigma_\mu\b{\sigma}_\nu + \sigma_\nu\b{\sigma}_\mu = 2\eta_{\mu\nu}\IM$;
    \item $\tr{(\sigma^\mu\b{\sigma}_\nu)} = 2\delta^\mu_\nu$.
\end{enumerate}
La corrispondenza tra quadrivettori e matrici hermitiane è invertibile. Infatti, le componenti del quadrivettore si possono ricavare dalla matrice $X$ come 
\begin{equation}
\label{eq: corrispondenza quadrivettori e matrici hermitiane}
    x^\mu = \frac{1}{2}\tr{(\sigma^\mu X)},\quad X := \b{\sigma}_\nu x^\nu.
\end{equation}
Si consideri quindi una matrice $A \in \group{SL}{2,\bb{C}}$. Con $A$ si può definire la trasformazione 
\begin{equation}
    X \to X' = AXA^\dag : X'^\dag = X^\dag,\,\det{X} = \det{X'}.
\end{equation}
Questa trasformazione induce una trasformazione lineare delle componenti del quadrivettore $x^{\mu\prime} \equiv \phi(A)^\mu_\nu x^\nu$, da cui, usando la proprietà \eqref{eq: corrispondenza quadrivettori e matrici hermitiane}, si trova che 
\begin{equation}
    \phi(A)^\mu_\nu = \frac{1}{2}\tr{(\sigma^\mu A\b{\sigma}_\nu A^\dag)}.
\end{equation}
Dato che la trasformazione $X \to AXA^\dag$ preserva la norma di Minkowski $x^\mu \eta_{\mu\nu} x^\nu$, significa che è una trasformazione di Lorentz, per cui 
\begin{eqnarray}
\label{eq: omomorfismo tra SL(2,C) e SO(1,3)}
    (\Lambda^{-1})^\mu_\nu = \phi(A)^\mu_\nu = \frac{1}{2}\tr{(\sigma^\mu A\bb{\sigma}_\nu A^\dag)}.
\end{eqnarray} 
In particolare, la mappa $\phi : A \to A' \equiv \phi(A)$ definisce un omomorfismo 
\begin{equation}
    \phi : \group{SL}{2,\bb{C}} \to \group{SO}{1,3}.
\end{equation}
Il nucleo dell'omomorfismo in \eqref{eq: omomorfismo tra SL(2,C) e SO(1,3)} è $\ker{\phi} = \{\IM, -\IM\}$, poiché entrambe inducono la stessa trasformazione. Chiaramente, $\ker{\phi} \cong \bb{Z}_2$ e, per il primo teorema di isomorfismo per i gruppi, si trova l'isomorfismo
\begin{equation}
    \group{SO}{1,3} \cong \group{SL}{2,\bb{C}}/\bb{Z}_2.
\end{equation}
Per questo, si dice che $\group{SL}{2,\bb{C}}$ è il \emph{doppio ricoprimento} del gruppo di Lorentz (ortocrono) $\group{SO}{1,3}$. Il gruppo di ricoprimento viene indicato tramite la nozione di \emph{gruppo di spin}, per cui $\group{Spin}{1,3} \cong \group{SL}{2,\bb{C}}$. Più in generale, in $D$ dimensioni, il gruppo di spin $\group{Spin}{1,D - 1}$ è il ricoprimento di $\group{SO}{1,D - 1}$ se la segnatura è quella di Minkowski, mentre $\group{Spin}{D}$ è il ricoprimento di $\group{SO}{D}$ se la segnatura è quella di Euclide. 

Si introducono ora due matrici che si possono dimostrare essere i generatori dell'algebra di Lorentz $\alg{so}{1,3}$ nelle rappresentazioni $\bm{(1/2,0)}$ e $\bm{(0,1/2)}$ rispettivamente. Tali matrici sono
\begin{equation}
    \sigma_{\mu\nu} = \dfrac{1}{4}(\sigma_\mu\b{\sigma}_\nu - \sigma_\nu\b{\sigma}_\mu),
\end{equation}
\begin{equation}
    \b{\sigma}_{\mu\nu} = \dfrac{1}{4}(\b{\sigma}_\mu\sigma_\nu - \b{\sigma}_\nu\sigma_\mu).
\end{equation}
Con queste si possono quindi scrivere le trasformazioni di Lorentz finite come
\begin{equation}
\label{eq: trasformazione di Lorentz finita L}
    L(\lambda) = e^{\lambda^{\mu\nu}\sigma_{\mu\nu}/2},
\end{equation}
\begin{equation}
\label{eq: trasformazione di Lorentz finita L tilde}
    \t{L}(\lambda) = e^{\lambda^{\mu\nu}\b{\sigma}_{\mu\nu}/2},
\end{equation}
dove $\lambda^{\mu\nu} = -\lambda^{\nu\mu}$ e, dato che $\sigma_{\mu\nu}^\dag = -\b{\sigma}_{\mu\nu}$, sono tali che
\begin{equation}
    L(\lambda)^\dag = \tilde{L}(\lambda)^{-1}.
\end{equation}
Le rappresentazioni matriciali $L(\lambda)$ ed $\t{L}(\lambda)$ di \eqref{eq: trasformazione di Lorentz finita L} ed \eqref{eq: trasformazione di Lorentz finita L tilde} sono direttamente in relazione con la mappa $X \to X' = AXA^\dag$. Per una trasformazione di Lorentz $\Lambda^{-1}$ legata ad $A \in \group{SL}{2,\bb{C}}$ dalla mappa \eqref{eq: omomorfismo tra SL(2,C) e SO(1,3)}, si può identificare $A \equiv L^{-1}$ ed $A^\dag \equiv \t{L}$. Questo si può dimostrare usando le relazioni $A\b{\sigma}_\mu A^\dag = (\Lambda^{-1})^\nu_\mu\b{\sigma}_\nu$ e $A^\dag \sigma_\mu A = (\Lambda^{-1})^\nu_\mu\sigma_\nu$.

Questa costruzione permette di descrivere campi che trasformano secondo le due rappresentazioni considerate inizialmente. Si considerino quindi due campi $\Phi(x)$ e $\chi(x)$, appartenenti rispettivamente alle rappresentazioni $\bm{(1/2,0)}$ e $\bm{(0,1/2)}$. Sotto una trasformazione di Lorentz $x^\mu \to x^{\prime\mu} \equiv \Lambda^\mu_\nu x^\nu$, le loro leggi di trasformazioni si possono scrivere come 
\begin{equation}
    \Psi(x) \to \Psi'(x') \equiv L^{-1}(\lambda)\Psi(\Lambda x),
\end{equation}
\begin{equation}
    \chi(x) \to \chi'(x') \equiv \tilde{L}^{-1}(\lambda)\chi(\Lambda x).
\end{equation}
In generale, i campi come $\Psi$ e $\chi$ prendono il nome di \emph{spinori di Weyl}, o \emph{spinori chirali}. In particolare, gli spinori che trasformano come $\Phi(x)$ sono detti \emph{spinori left}, in quanto trasformano secondo la \emph{rappresentazione sinistrorsa} $\bm{(1/2,0)}$, mentre quelli che trasformano come $\chi(x)$ sono detti \emph{spinori right}, in quanto trasformano con la \emph{rappresentazione destrorsa} $\bm{(0,1/2)}$.

\section{Equazione di Dirac}
I campi spinoriali furono introdotti da Dirac per descrivere i fermioni. Questi campi soddisfano un'equazione del moto, che è l'equazione di Dirac
\begin{equation}
\label{eq: equazione di Dirac}
    (\sl{\p} - m)\Psi(x) = 0,\quad \sl{\p} \equiv \gamma^\mu\p_\mu,
\end{equation}
dove, contestualmente, si sono introdotte le matrici quadrate $\gamma^\mu$, agenti sugli indici dello spinore $\Phi(x)$, che prendono il nome di \emph{matrici di Dirac}. In particolari, queste matrici soddisfano la relazione di anti-commutazione 
\begin{equation}
\label{eq: algebra di Clifford}
    \acomm{\gamma^\mu}{\gamma^\nu} \equiv \gamma^\mu\gamma^\nu + \gamma^\nu\gamma^\mu = 2\eta^{\mu\nu}\IM_D.
\end{equation}
La condizione \eqref{eq: algebra di Clifford} definisce l'\emph{algebra di Clifford} per uno spaziotempo $D$-dimensionale, indicata con $\alg{Cliff}{D}$. Le $D$ matrici $\gamma^\mu$ sono i generatori dell'algebra di Clifford, una cui base consiste nella matrice identità insieme a tutti i prodotti indipendenti dei generatori. Per l'algebra di Clifford esiste una rappresentazione tramite matrici quadrate $2^{[D/2]}$-dimensionali, dove $[x]$ rappresenta l'intero più grande minore od uguale ad $x$. La rappresentazione è unica a meno di un'equivalenza per $D = 2m$, mentre per $D = 2m + 1$ le rappresentazioni inequivalenti sono due. Inoltre, è sempre possibile scegliere una rappresentazione dove le matrici $\gamma^i$ sono hermitiane, mentre $\gamma^0$ è anti-hermitiana. 


In generale, spinori e matrici di Dirac si possono definire per ogni dimensione $D$; in particolare, le $\gamma^\mu$ in dimensione $D$ generica sono necessariamente complesse. Qui ci si concentra specialmente sul caso $D = 4$. Rappresentazioni equivalenti dell'algebra di Clifford descrivono la stessa fisica. Pertanto, le implicazioni fisiche dell'equazione del formalismo di Dirac deve risultare indipendente dalla scelta della rappresentazione. Effettivamente, gran parte della fisica può essere dedotta in una maniera indipendente dalla rappresentazione, ma una rappresentazione esplicita è conveniente per alcuni scopi. A tal proposito, si introduce la \emph{rappresentazione di Weyl} in $D = 4$ dove le matrici di Dirac $\gamma^\mu$ sono scritte in funzione delle matrici di Weyl sull'anti-diagonale, viste in equazione \eqref{eq: matrici di Weyl}, ossia 
\begin{equation}
\label{eq: rappresentazione di Weyl per le matrici di Dirac}
    \gamma^\mu := 
    \begin{pmatrix}
        \OM_2 & \sigma^\mu \\
        \b{\sigma}^\mu & \OM_2
    \end{pmatrix},\quad \dim{\gamma^\mu} = 4.
\end{equation}
Ci sono rappresentazioni ``anti-diagonali a blocchi'' di questo tipo per ogni dimensione $D$ pari. 

Un fatto fondamentale riguardo all'algebra di Clifford in dimensione $D$ è che i commutatori 
\begin{equation}
    \Sigma^{\mu\nu} := \frac{1}{4}\comm{\gamma^\mu}{\gamma^\nu},
\end{equation}
sono i generatori di una rappresentazione $2^{[D/2]}$-dimensionale dell'algebra di Lie $\alg{so}{1, D - 1}$. Nella rappresentazione di Weyl, si può vedere che le $\Sigma^{\mu\nu}$ sono diagonali a blocchi con le matrici $\sigma^{\mu\nu}$ e $\b{\sigma}^{\mu\nu}$ come entrate, cioè
\begin{equation}
    \Sigma^{\mu\nu} = 
    \begin{pmatrix}
        \sigma^{\mu\nu} & \OM_2 \\
        \OM_2 & \b{\sigma}^{\mu\nu}
    \end{pmatrix}.
\end{equation}
La rappresentazione 4-dimensionale dell'algebra $\alg{so}{1, 3}$ è quindi riducibile. In particolare, è la somma diretta delle rappresentazioni $\bm{(1/2,0)}$ e $\bm{(0,1/2)}$ viste prima. Chiaramente, la riducibilità di $\Sigma^{\mu\nu}$ vale per qualsiasi scelta delle $\gamma^\mu$, dato che le matrici $\Sigma^{\mu\nu}$ sono equivalenti a quelle della rappresentazione di Weyl. In tutte le dimensioni pari, esiste una riduzione analoga alla somma diretta di una rappresentazione irriducibile $2^{[D/2] - 1}$-dimesionale di $\alg{so}{1, 3}$ insieme alla rappresentazione coniugata. In dimensione dispari, la rappresentazione data dalle $\Sigma^{\mu\nu}$ è irriducibile. Le trasformazioni proprie e finite di Lorentz sono rappresentate dalle matrici 
\begin{equation}
\label{eq: matrici trasformazioni di Lorentz finite proprie}
    L(\lambda) := e^{\lambda^{\mu\nu}\Sigma_{\mu\nu}} \in \group{Spin}{1,D - 1},
\end{equation}
e si ha che 
\begin{equation}
    L(\lambda)^{-1}\gamma^\mu L(\lambda) = \Lambda(\lambda)^\mu_\nu \gamma^\nu,
\end{equation}
che corrisponde esattamente ad una trasformazione di Lorentz
\begin{equation}
    \gamma^\mu \to \gamma^{\prime\mu} \equiv \Lambda(\lambda)^\mu_\nu \gamma^\nu = L(\lambda)^{-1}\gamma^\mu L(\lambda).
\end{equation}
Da questo si può dedurre l'importante proprietà di covarianza di Lorentz dell'equazione di Dirac: data qualsiasi soluzione $\Psi(x)$ anche 
\begin{equation}
    \Psi(x) \to \Psi'(x') = L(\lambda)^{-1}\Psi(\Lambda(\lambda) x)
\end{equation}
è una soluzione. 
\begin{ex}[rotazione di uno spinore]
    Considerando la particolare trasformazione $\lambda^{12} \equiv \theta \neq 0$, si ha 
    \begin{equation}
        L(\lambda) = e^{\theta\Sigma_{12}} = 1 + \theta\Sigma_{12},
    \end{equation}
    dove le potenze pari sono proporzionali all'identità; infatti
    \begin{equation}
        \Sigma_{12}\Sigma_{12} = \frac{1}{2}(\gamma^1\gamma^2)\frac{1}{2}(\gamma^1\gamma^2) = -\frac{1}{4}\IM_4.
    \end{equation}
    Quindi, si ottiene la relazione
    \begin{equation}
        \cos{\bigg(\frac{\theta}{2}\bigg)}\IM + 2\sin{\bigg(\frac{\theta}{2}\bigg)}\Sigma_{12},
    \end{equation}
    da cui si vede che per $\theta = 2\pi$, si ottiene $-\IM$. 
\end{ex}
Per ricavare l'equazione \eqref{eq: equazione di Dirac}, si deve definire un'azione per il campo spinoriale. La sua scelta è limitata dalle condizioni:
\begin{enumerate}
    \item l'azione deve essere reale;
    \item l'azione deve essere lineare nelle derivate;
    \item l'azione deve essere un invariante di Lorentz.
\end{enumerate}
Un termine invariante lineare nelle derivate è $\gamma^\mu\p_\mu$. Per quanto riguarda il campo $\Psi$, questo deve comparire nell'espressione dell'azione attraverso un invariante di Lorentz. Per un primo tentativo, si nota che su $\bb{C}^4$ esiste un prodotto interno naturale dato da
\begin{equation}
    \Psi^\dag\Psi \in \bb{R}.
\end{equation}
Essendo uno scalare, si potrebbe pensare che $\Psi^\dag\Psi$ sia anche invariante di Lorentz; tuttavia, usando le trasformazioni
\begin{equation}
    \begin{cases}
        \Psi \to \Psi' \equiv L^{-1}(\lambda)\Psi(\Lambda x) \\
        \Psi^\dag \to \Psi'^\dag \equiv \Psi^\dag(\Lambda x)[L^{-1}(\lambda)]^\dag
    \end{cases},
\end{equation}
si trova che $\Psi^\dag\Psi$ non è in effetti un invariante di Lorentz. Si può quindi assumere che l'invariante di Lorentz cercato sia un bilineare non-degenere della forma 
\begin{equation}
    \Psi^\dag\beta\Psi,
\end{equation}
dove $\beta$ è una matrice quadrata ed invertibile da determinare. L'invarianza di Lorentz richiede che 
\begin{equation}
\label{eq: condizione di L-invarianza su beta}
    \Sigma_{\mu\nu}^\dag\beta + \beta\Sigma_{\mu\nu} = 0.
\end{equation}
Dato che si sta cercando una forma bilineare reale, si impone che $\beta$ sia una matrice hermitiana. Per quanto detto prima, $\beta$ sicuramente non può essere scelta come l'identità, anche perché, se così fosse, ciò implicherebbe che il gruppo di Lorentz sia compatto, quando in realtà è non-compatto. Tuttavia, si può verificare che la condizione \eqref{eq: condizione di L-invarianza su beta} è validata se si sceglie $\beta$ essere un qualsiasi multiplo di $\gamma^0$. Questa forma è valida in qualsiasi rappresentazione hermitiana ed è unica. Quindi, si hanno
\begin{equation}
    \beta\gamma_\mu\beta^{-1} = -(\gamma_\mu)^\dag,\quad \beta\Sigma_{\mu\nu}\beta^{-1} = -(\Sigma_{\mu\nu})^\dag.
\end{equation}
Tipicamente, si opera la scelta 
\begin{equation}
    \beta := i\gamma^0.
\end{equation}
Di conseguenza, si definisce l'\emph{aggiunto di Dirac} come 
\begin{equation}
\label{eq: aggiunto di Dirac}
    \b{\Psi} := \Psi^\dag \beta = i\Psi \gamma^0.
\end{equation}
La forma bilineare invariante cercata si può dunque scrivere come 
\begin{equation}
    \b{\Psi}\Psi.
\end{equation}
Si può finalmente definire l'azione per il campo di Dirac libero. Questa è
\begin{equation}
    S[\Psi,\b{\Psi}] = \int d^4x\,(m\b{\Psi}\Psi - \b{\Psi}\sl{\p}\Psi).
\end{equation}
Dal principio di minima azione applicato alla variazione 
\begin{equation}
    \delta S = \int d^4x\,[\delta\b{\Psi}(\sl{\p} - m)\Psi - (\p_\mu\b{\Psi}\gamma^\mu + \b{\Psi}m)\delta\Psi],
\end{equation}
si ottengono le due equazioni del moto equivalenti
\begin{equation}
    \begin{cases}
        (\sl{\p} - m)\Psi = 0 \\
        \b{\Psi}(\backv{\sl{\p}} - m) = 0
    \end{cases},\quad \backv{\sl{\p}} \equiv \p_\mu\b{\Psi}\gamma^\mu.
\end{equation}

\subsection{Spinori di Weyl in dimensione pari}
Tornando al caso $D = 2m$, dato che esiste una rappresentazione anti-diagonale a blocchi per ogni dimensione pari, lo stesso vale per tutte le dimensioni $D = 2m$ con due sotto-rappresentazioni irriducibili, ognuna di dimensione $2^{m - 1}$. Questo suggerisce che uno spinore di Dirac $\Psi(x)$ non è il campo covariante di Lorentz più semplice. In un certo senso questo è vero: si può definire uno spinore di Weyl $\psi(x)$ con $2^{m - 1}$ componenti che trasforma secondo 
\begin{equation}
    \psi(x) \to \psi'(x') \equiv L(\lambda)^{-1}\psi(\Lambda(\lambda)x),
\end{equation}
dove le $L(\lambda)$ sono definite come in \eqref{eq: matrici trasformazioni di Lorentz finite proprie}, ma ora per ogni dimensione pari. Si può anche definire un campo $\b{\chi}$ che trasforma secondo la rappresentazione coniugata, cioè come 
\begin{equation}
    \b{\chi} \to \b{\chi}'(x') \equiv \t{L}(\lambda)^{-1}\b{\chi}(\Lambda(\lambda)x),
\end{equation}
dove le $L(\lambda)$ sono definite di nuovo come in \eqref{eq: matrici trasformazioni di Lorentz finite proprie}. Per questi campi, valgono le \emph{equazioni di Weyl}
\begin{align}
\label{eq: equazioni di Weyl}
    \b{\sigma}^{\mu\nu}\p_\mu\psi(x) &= 0, \\
    \sigma^{\mu\nu}\p_\mu\b{\chi}(x) &= 0,
\end{align}
le quali, ammettendo soluzioni ad onda piana, descrivono particelle prive di massa, cioè ``massless''. Per scrivere un'azione per lo spinore di Weyl serve considerare $\psi(x)$ insieme al suo aggiunto $\psi(x)^\dag$, con i quali si può costruire una forma bilineare che trasforma come un vettore. L'azione per $\psi(x)$ è allora 
\begin{equation}
    S[\psi,\b{\psi}] = -\int d^Dx\,i\psi^\dag\b{\sigma}^\mu\p_\mu\psi,
\end{equation}
di cui ne esiste una versione analoga per $\b{\chi}(x)$. Si noti infine che non c'è modo di introdurre la massa usando un singolo spinore di Weyl; infatti, l'equazione 
\begin{equation}
    \b{\sigma}^\mu\p_\mu\psi(x) = m\psi(x)
\end{equation}
non è Lorentz-invariante. Invece, per descrivere particelle massive serve usare insieme $\psi(x)$ e $\b{\chi}$. Infatti, si può vedere che uno spinore di Dirac può essere scritto come 
\begin{equation}
    \Psi(x) = 
    \begin{pmatrix}
        \psi(x) \\
        \b{\chi}(x)
    \end{pmatrix},
\end{equation}
il quale, se inserito nell'equazione di Dirac \eqref{eq: equazione di Dirac} nella rappresentazione di Weyl \eqref{eq: rappresentazione di Weyl per le matrici di Dirac}, permette di dimostrare che l'equazione \eqref{eq: equazione di Dirac} è equivalente alla coppia di equazioni 
\begin{equation}
    \begin{cases}
        \b{\sigma}^\mu\p_\mu\psi(x) = m\b{\chi}(x) \\
        \sigma^\mu\p_\mu\b{\chi}(x) = m\psi(x)
    \end{cases}.
\end{equation}

\section{Algebra di Clifford}
L'obiettivo ora è discutere l'algebra di Clifford associata al gruppo di Lorentz in $D$ dimensioni. Si parte dalla costruzione esplicita e generale delle matrici $\gamma$ euclidee, le quali soddisfano la condizione \eqref{eq: algebra di Clifford} con $\eta_{\mu\nu}$ sostituita da $\delta_{\mu\nu}$. Quindi, si hanno
\begin{gather}
\label{eq: costruzione matrici di Dirac euclidee}
    \gamma^1 = \sigma_1 \ot \IM \ot \cdots \ot \IM, \\
    \gamma^2 = \sigma_2 \ot \IM \ot \cdots \ot \IM, \\
    \gamma^3 = \sigma_3 \ot \sigma_1 \ot \IM \ot \cdots \ot \IM, \\
    \gamma^4 = \sigma_3 \ot \sigma_2 \ot \IM \ot \cdots \ot \IM, \\
    \gamma^5 = \sigma_3 \ot \sigma_3 \ot \sigma_1 \ot \IM \ot \cdots \ot \IM, \\
    \gamma^6 = \sigma_3 \ot \sigma_3 \ot \sigma_2 \ot \IM \ot \cdots \ot \IM, 
\end{gather}
e così via, dove si è usato il prodotto tensoriale, o di Kronecker, definito da
\begin{equation}
    A \ot B = 
    \begin{pmatrix}
        A_{11}B & \cdots & A_{1n}B \\
        \vdots & \ddots & \vdots \\
        A_{n1}B & \cdots & A_{nn}B
    \end{pmatrix},\quad \dim{(A \ot B)} = nm \times nm.
\end{equation}
Tutte queste matrici sono hermitiane, con il cui quadrato pari all'identità $\IM$ e mutuamente anti-commutanti. Assumendo $D = 2m$, sono necessari $m$ fattori in questa costruzione per ottenere $\gamma^\mu$ con $1 \le \mu \le D = 2m$. Pertanto, si ottiene una rappresentazione $2^{D/2}$-dimensionale. Dall'altra parte, nel caso in cui $D = 2m + 1$, è necessaria un'ulteriore matrice. Si sceglie allora $\gamma^{2m + 1}$ dalla lista di matrici elencata sopra, ma mantenendo solamente i primi $m$ fattori, in altre parole, eliminando una $\sigma_1$. In questo modo, la dimensione della rappresentazione non aumenta passando da $D = 2m$ a $D = 2m + 1$. Si può allora dire che la costruzione in \eqref{eq: costruzione matrici di Dirac euclidee} restituisce una rappresentazione di dimensione $2^{[D/2]}$. 

Anche se le matrici $\gamma$ euclidee hanno applicazioni fisiche, si è qui interessati alla loro controparte lorentziana. Per ottenerle, si deve semplicemente scegliere una qualsiasi matrice tra quelle euclidee, moltiplicarla per $i$ e indicarla con $\gamma^0$ per la direzione ``time-like''. Questa matrice è anti-hermitiana e soddisfa $(\gamma^0)^2 = -\IM$. Conseguentemente, si rinominano le rimanenti $D - 1$ matrici così da ottenere l'insieme $\gamma^\mu$ con $\mu = 0,\dots, D - 1$. Le proprietà di hermiticità delle matrici $\gamma^\mu$ sono riassunte nella relazione 
\begin{equation}
\label{eq: hermiticità matrici di Dirac}
    (\gamma^\mu)^\dag = \gamma^0\gamma^\mu\gamma^0.
\end{equation}
\'E fondamentale per la teoria di Dirac che la fisica degli spinori sia la medesima in ogni rappresentazione equivalente dell'algebra di Clifford. Quindi, si è di fatto interessati alle classi di rappresentazioni legate dalla coniugazione, cioè
\begin{equation}
\label{eq: matrici gamma legate da coniugazione}
    \gamma^{\prime\mu} = S\gamma^\mu S^{-1},
\end{equation}
dove le matrici $S$, dato che si considerano solo rappresentazioni hermitiane, devono essere unitarie. Date due rappresentazioni equivalenti, la matrice di trasformazione $S$ è unica a meno di un fattore di fase. A meno di questa coniugazione, esiste un'unica rappresentazione irriducibile dell'algebra di Clifford tramite matrici $2^m \times 2^m$ per $D = 2m$. Ogni altra rappresentazione è riducibile ed equivalente alla somma diretta di due copie della rappresentazione irriducibile più sopra. Si può sempre scegliere una rappresentazione irriducibile hermitiana, definita in modo da soddisfare \eqref{eq: hermiticità matrici di Dirac}. In dimensione dispari, esistono due rappresentazioni irriducibili inequivalenti che differiscono nel segno della matrice ``finale'' $\pm \gamma^{2m + 1}$. Le conseguenze fisiche sono indipendenti dalla particolare scelta della rappresentazione. 
\begin{ex}[matrici di Dirac in $D = 2,3,4$]
    In $D = 2$, si devono trovare due matrici $\gamma^1$ e $\gamma^2$ tali che $\acomm{\gamma^1}{\gamma^2} = 0$ e che $(\gamma^1)^2 = (\gamma^2)^2 = \IM_4$. Una scelta possibile è $\gamma^1 = \sigma_1$ e $\gamma^2 = \sigma_2$, definendo $\alg{Cliff}{2}$. Per $D = 3$ si deve aggiungere $\gamma^3 = \sigma_3$, definendo così $\alg{Cliff}{3}$. In $D = 4$ si scelgono invece
    \begin{equation}
        \begin{cases}
            \gamma^0 = \sigma_1 \ot \IM \\
            \gamma^1 = \sigma_2 \ot \IM \\
            \gamma^2 = \sigma_3 \ot \sigma_1 \\
            \gamma^3 = \sigma_3 \ot \sigma_2
        \end{cases},
    \end{equation}
    ottenendo $\alg{Cliff}{4}$. 
\end{ex}
% Si torni in $D = 4$. Qui si ha 
% \begin{equation}
%     \Sigma^{\mu\nu} = 
%     \begin{pmatrix}
%         \sigma_{\mu\nu} & \OM_2 \\
%         \OM_2 & \b{\sigma}_{\mu\nu}
%     \end{pmatrix},
% \end{equation}
% andando ad agire su $\Psi(x)$ le due componenti sono indipendenti,
% \begin{equation}
%     \Psi(x) = 
%     \begin{pmatrix}
%         \Psi(x) \\
%         \chi(x)
%     \end{pmatrix},
% \end{equation}
% dove la componente left trasforma con $\sigma_{\mu\nu}$, mentre quella right con $\b{\sigma}_{\mu\nu}$. Tutto ciò è ripetibile in $D$ dimensioni, ottenendo la relazione
% \begin{equation}
%     \underbrace{2^{D/2}}_{rappr.\,riducibile} = \underbrace{2^{D/2 - 1}}_{irrep} + \underbrace{2^{D/2 - 1}}_{irrep}.
% \end{equation}

\subsection{Algebra di Clifford completa}
L'algebra di Clifford completa consiste nella matrice identità $\IM$, i $D$ generatori $\gamma$ e tutte le matrici indipendenti costruite come prodotti dei generatori. Dato che i prodotti simmetrici si riducono in prodotti contenenti un numero minore di matrici $\gamma$ a causa della condizione \eqref{eq: algebra di Clifford}, i nuovi elementi devono essere prodotti anti-simmetrici. Si definisce quindi 
\begin{equation}
\label{eq: anti-simmetrizzazione prodotti matrici gamma}
    \gamma^{\mu_1\cdots\mu_r} = \gamma^{[\mu_1}\gamma^{\mu_2}\cdots\gamma^{\mu_r]}.
\end{equation}
Componenti non nulle di tensori si possono quindi scrivere come 
\begin{equation}
    \gamma^{\mu_1\cdots\mu_r} = \gamma^{\mu_1}\cdots\gamma^{\mu_r},\quad \mu_1 \neq \cdots \neq \mu_r.
\end{equation}
Tutte queste matrici sono a traccia nulle, a parte al rango minimo per $r = 0$, che corrisponde alla matrice identità, e al rango massimo per $r = D$, che è a traccia nulla solo per dimensioni $D$ pari. 

Le teorie di supersimmetria e la supergravità emergono dal concetto di spin fermionico. Come ci si può aspettare, sono richieste proprietà intricate dell'algebra di Clifford per determinare le conseguenze fisiche di queste teorie di campo. Si riportano perciò nel seguito alcuni risultati utili riguardo alle matrici $\gamma$, validi sia in dimensione parti che in dimensione dispari. Si parte dalle possibili contrazioni delle matrici $\gamma$, due primi esempi sono 
\begin{equation}
    \gamma^\mu\gamma_\mu = D,\quad \gamma^{\mu\nu}\gamma_\nu = (D - 1)\gamma^\mu.
\end{equation}
La logica dietro seconda relazione, e quindi anche dietro alla prima, sta semplicemente nel fatto che l'indice $\nu$ assume tutti i valori eccetto $\mu$, risultando in $D - 1$ termini nella somma. Una logica simile spiega anche 
\begin{equation}
    \gamma^{\mu\nu\rho}\gamma_\rho = (D - 2)\gamma^{\mu\nu}.
\end{equation}
In generale, si dimostra 
\begin{equation}
\label{eq: contrazione generale matrici di Dirac}
    \gamma^{\mu_1\cdots\mu_r\nu_1\cdots\nu_s}\gamma_{\nu_s\cdots \nu_1} = \frac{(D - r)!}{(D - r - s)!}\gamma^{\mu_1\cdots\mu_r}.
\end{equation}
Si noti che la seconda matrice a primo membro ha gli indice in ordine inverso, cosicché non compaia un segno meno. In generale, quindi, è molto utile conoscere la regola di inversione
\begin{equation}
\label{eq: regola inversione indici matrici di Dirac}
    \gamma^{\nu_1\cdots\nu_r} = (-1)^{r(r - 1)/2}\gamma^{\nu_r\cdots\nu_1},
\end{equation}
dove il fattore a secondo membro è negativo per $r = 2,3 \mod 4$. 

Anche senza contrarre gli indici, esistono alcune relazioni utili. Ad esempio, quando si deve calcolare un prodotto del tipo $\gamma^{\mu_1\mu_2}\gamma_{\nu_1\cdots\nu_D}$, si sa che i valori degli indici $\mu_1$ e $\mu_2$ compaiono nell'insieme di indici $\{\nu_i\}$. Ci sono $D$ possibilità per $\mu_2$, e dato che $\mu_1$ deve essere diverso, rimangono $D - 1$ possibilità per $\mu_1$. Di conseguenza, il risultato è 
\begin{equation}
    \gamma^{\mu_1\mu_2}\gamma_{\nu_1\cdots\nu_D} = D(D - 1)\delta^{\mu_1\mu_2}_{[\nu_1\nu_2}\gamma_{\nu_3\cdots\nu_D]},
\end{equation}
da cui si ottiene anche 
\begin{equation}
    \gamma^{\mu_1\mu_2}\gamma_{\nu_1\cdots\nu_D}\epsilon^{\nu_1\cdots\nu_D} = D(D - 1)\epsilon^{\mu_2\mu_1\nu_3\cdots\nu_D}\gamma_{\nu_3\cdots\nu_D}.
\end{equation}

Considerando il prodotto di matrici $\gamma$ senza contrazione degli indici, il caso più semplice è $\gamma^\mu\gamma^\nu$, il quale si può dimostrare essere 
\begin{equation}
    \begin{split}
        \gamma^\mu\gamma^\nu &= \frac{1}{2}(\gamma^\mu\gamma^\nu + \gamma^\nu\gamma^\mu) + \frac{1}{2}(\gamma^\mu\gamma^\nu - \gamma^\nu\gamma^\mu) \\
        &= \eta^{\mu\nu} + \frac{1}{2}(\gamma^\mu\gamma^\nu - \gamma^\nu\gamma^\mu) \\
        &\equiv \eta^{\mu\nu} + \gamma^{[\mu}\gamma^{\nu]} \\
        &\equiv \eta^{\mu\nu} + \gamma^{\mu\nu}.
    \end{split}
\end{equation}
Il risultato discende direttamente dalle definizioni: la parte antisimmetrica del prodotto è definita in \eqref{eq: anti-simmetrizzazione prodotti matrici gamma}, in particolare da 
\begin{equation}
    \gamma^{[\mu}\gamma^{\nu]} \equiv \gamma^{\mu\nu} \equiv \frac{1}{2}(\gamma^\mu\gamma^\nu - \gamma^\nu\gamma^\mu),
\end{equation}
mentre la parte simmetrica del prodotto è $\eta^{\mu\nu}$ per la condizione \eqref{eq: algebra di Clifford}. Questo illustra l'approccio generale: si scrive prima la matrice di Clifford totalmente anti-simmetrica che contiene tutti gli indici e poi si aggiungono i termini per tutte le possibili gli accoppiamenti degli indici. Un altro esempio è
\begin{equation}
    \gamma^{\mu\nu\rho}\gamma_{\sigma\tau} = \gamma^{\mu\nu\rho}_{\sigma\tau} + 6\gamma^{[\mu\nu}_{[tau}\delta^{\rho]}_{\sigma]} + 6\gamma^{[\mu}\delta^{\nu}_{[\tau}\delta^{\rho]}_{\sigma]}.
\end{equation}
Considerando infine prodotti di matrici dove sono presenti indici contratti e non contratti, si può dimostrare in generale che 
\begin{equation}
    \gamma_\rho\gamma^{\mu_1\cdots\mu_r}\gamma^\rho = (-1)^r(D - 2r)\gamma^{\mu_1\cdots\mu_r}.
\end{equation}
% In dimensioni pari esiste una sola rappresentazione irriducibile, mentre in dimensione dispari ne esistono due ma portano alla stessa algebra. L'algebra contiene $\{\IM,\gamma^\mu,\gamma^\mu\gamma^\nu,\gamma^\mu\gamma^\nu\gamma^\rho\}$, ma in realtà la definizione di algebra di Clifford permette di restringersi ad un insieme più piccolo.

\subsubsection{Base dell'algebra per \texorpdfstring{$D$}{D} pari}
Ci si concentra ora sullo spaziotempo di dimensione $D$ pari e si costruisce una base ortogonale per l'algebra di Clifford. Successivamente, si potranno estendere i risultati per le dimensioni dispari. La base si indica con $\{\Gamma^A\}$ ed è
\begin{equation}
    \{\Gamma^A\} \equiv \{\IM,\gamma^\mu,\gamma^{\mu_1\mu_2},\gamma^{\mu_1\mu_2\mu_3},\dots,\gamma^{\mu_1\cdots\mu_D}\},
\end{equation}
dove i valori degli indici soddisfano la condizione $\mu_1 < \cdots < \mu_r$. Può essere utile definire anche la base di ordine inverso 
\begin{equation}
    \{\Gamma_A\} \equiv \{\IM,\gamma_\mu,\gamma_{\mu_2\mu_1},\gamma_{\mu_3\mu_2\mu_1},\dots,\mu_{\mu_D\cdots\mu_1}\}.
\end{equation}
Per \eqref{eq: regola inversione indici matrici di Dirac}, le matrici in $\{\Gamma^A\}$ e $\{\Gamma_A\}$ differiscono al più per un fattore di segno. La lista di matrici $\{\Gamma^A\}$ contiene $2^D$ matrici ortogonali in un'algebra di dimensione totale $2^D$. Pertanto, risulta essere una base dello spazio di matrici $M$ di dimensione $2^m \times 2^m$. Si può dimostrare poi che $\Gamma^A\Gamma^B = \pm \Gamma^C$, dove $\Gamma^C$ è un elemento della base i cui indici sono comuni a quelli di $A$ e $B$ con quelli comuni esclusi, e che 
\begin{equation}
    \tr{(\Gamma^A\Gamma_B)} = 2^m\delta^A_B.
\end{equation}
La conclusione che esiste un'unica rappresentazione dell'algebra di Clifford di dimensione $2^m$ è vera per ogni metrica pseudo-euclidea di segnatura $(p,q)$. Un altro fatto generale è che gli elementi di rango due $\gamma^{\mu\nu}$ dell'algebra di Clifford sono i generatori di una rappresentazione dell'algebra di Lie $\alg{so}{p, q}$ con $ p + q = D = 2m$ con una metrica di segnatura $(p, q)$. 

\subsubsection{Matrice di chiralità}
Per molte ragioni può essere utile studiare l'elemento tensoriale di rango più alto dell'algebra di Clifford, dato che fornisce un legame tra dimensione pari e dimensione dispari ed è strettamente legato alla chiralità dei fermioni. Si definisce la matrice di chiralità
\begin{equation}
\label{eq: matrice di chiralità}
    \gamma_* := (-i)^{m + 1}\gamma^0\cdots \gamma^{D - 1},
\end{equation}
tale che $(\gamma_*)^2 = \IM$ per ogni dimensione pari, $\tr{(\gamma_*)} = n_+ - n_- = 0 = 0$, dove $n_+$ è il numero di autovalori positivi $+1$, mentre $n_-$ è quello di autovalori negativi $-1$ tali che $n_+ = n_- \equiv n$ e $n = 2^{m - 1}$, ed è hermitiana. Inoltre, si dimostra che $\acomm{\gamma_*}{\gamma^\mu} = 0$ e che $\comm{\gamma_*}{\gamma^{\mu\nu}} = 0$. Dato che $\gamma_*^2 = \IM$ e che $\tr{(\gamma_*)} = 0$, segue che si può scegliere una rappresentazione dove 
\begin{equation}
    \gamma_* = 
    \begin{pmatrix}
        \IM_{2^{m - 1}} & \OM \\
        \OM & -\IM_{2^{m - 1}}
    \end{pmatrix}.
\end{equation}
La matrice di chiralità risulta essere l'unico elemento di rango massimo nella lista $\{\Gamma^A\}$. Per ogni ordine delle componenti $\mu_i$, si può scrivere 
\begin{equation}
    \gamma_{\mu_0\cdots\mu_D} = i^{m + 1}\epsilon_{\mu_0\cdots\mu_D}\gamma_*.
\end{equation}
Tramite la matrice di chiralità di possono anche definire i \emph{proiettori di chiralità} con i quali si possono ottenere gli spinori di Weyl $\psi(x)$ e $\b{\chi}(x)$ a partire da uno spinore di Dirac $\Psi(x)$. Questi sono definiti come 
\begin{equation}
\label{eq: proiettori di chiralità}
    \begin{cases}
        P_L := \dfrac{1}{2}(\IM + \gamma_*) \\[2ex]
        P_R := \dfrac{1}{2}(\IM - \gamma_*)
    \end{cases},
\end{equation}
per cui 
\begin{equation}
    \begin{pmatrix}
        \psi(x) \\
        0
    \end{pmatrix}
    \equiv P_L\Psi(x),\quad 
    \begin{pmatrix}
        0 \\
        \b{\chi}(x)
    \end{pmatrix} 
    \equiv P_R\Psi(x).
\end{equation}
Queste relazioni ovviamente non sono vere in generale, ma valgono solo nella rappresentazione di Weyl. Si può dimostrare che le matrici in \eqref{eq: proiettori di chiralità} proiettano su sottospazi ortogonali, cioè che 
\begin{equation}
    P_L^2 = P_L,\quad P_R^2 = P_R,\quad P_LP_R = P_RP_L = 0,
\end{equation}
indipendentemente dalla rappresentazione dell'algebra di Clifford. 

\subsubsection{Base dell'algebra per \texorpdfstring{$D$}{D} dispari}
Per definire una base per l'algebra di Clifford in dimensione $D = 2m + 1$ è sufficiente riorganizzare le matrici nell'algebra di Clifford per $D = 2m$. In particolare, si definiscono due insiemi di $2m + 1$ generatori aggiungendo l'elemento di rango massimo $\gamma_*$ come segue 
\begin{equation}
\label{eq: basi per algebra di Clifford in D dispari}
    \gamma_\pm^\mu = \{\gamma^0,\gamma^1,\dots,\gamma^{2m - 1},\gamma^{2m} \equiv \pm\gamma_*\}.
\end{equation}
In questo modo si ottengono due insiemi di matrici, entrambi tali da soddisfare \eqref{eq: algebra di Clifford} per $D = 2m + 1$. I due insiemi $\{\gamma^\mu_+\}$ e $\{\gamma^\mu_-\}$ non sono equivalenti, ma portano a rappresentazioni equivalenti del gruppo di Lorentz. La differenza principale con il caso $D = 2m$ è che le matrici in $\{\Gamma^A\}$ non sono tutte indipendenti e formano quindi un insieme più che completo. In effetti, l'elemento di rango massimo di questa lista, che consiste nel prodotto di tutte le matrici $\gamma$, è, per \eqref{eq: basi per algebra di Clifford in D dispari}, un fattore di fase moltiplicato per la matrice identità. Più in generale, i settori di rango $r$ e di rango $D - r$ sono messi in relazione dalle relazioni di dualità 
\begin{equation}
    \gamma_\pm^{\mu_1\cdots\mu_r} = \pm\frac{i^{m + 1}}{(D - r)!}\epsilon^{\mu_1\cdots\mu_D}\gamma_{\pm\,\mu_D\cdots\mu_{r + 1}}.
\end{equation}
In dimensione pari si può trovare una relazione analoga
\begin{equation}
    \gamma^{\mu_1\cdots \mu_r}\gamma^* = -\frac{(-i)^{m + 1}}{(D - r)!}\epsilon^{\mu_r\cdots\mu_1\nu_1\cdots\nu_{D - r}}\gamma_{\nu_1\cdots \nu_{D - r}}.
\end{equation}
Pertanto, una base dell'algebra di Clifford in $D = 2m + 1$ dimensioni contiene le matrici in $\{\Gamma^A\}$ solo fino al rango $m$.

% Si definisce allora
% \begin{equation}
%     \gamma^\mu = 
%     \begin{pmatrix}
%         \alpha^\mu & \sigma^\mu \\
%         \b{\sigma}^\mu & \beta^\mu
%     \end{pmatrix},
% \end{equation}
% dove, per la relazione dell'algebra di Clifford, si ha $\alpha^\mu = \beta^\mu = 0$ e 
% \begin{equation}
%     \sigma^\mu\b{\sigma}^\nu + \sigma^\nu\b{\sigma}^\mu = \acomm{\sigma^\mu}{\b{\sigma}^\nu} = 2\eta^{\mu\nu},
% \end{equation}
% dove si definiscono le matrici 
% \begin{equation}
%     \sigma^\mu = (\IM,\sigma^i),\quad \b{\sigma}^\mu = (-\IM,\sigma^i),
% \end{equation}
% con $\sigma^i$ la generalizzazione delle matrici di Pauli, in quattro dimensioni $\acomm{\sigma^i}{\sigma^j} = \delta^{ij}$. Tutto ciò vale nella base di Weyl. 

\subsection{Simmetrie delle matrici \texorpdfstring{$\gamma$}{gamma}}
Ricordando che il coniugato di Dirac è definito come $\b{\Psi} = \Psi^\dag\gamma^0$, si definisce $A := i\gamma^0$, tale che
\begin{equation}
    A\gamma^\mu A^{-1} = -(\gamma^\mu)^\dag
    \quad\Longleftrightarrow\quad
    A\gamma^\mu=-(\gamma^\mu)^\dag A.
\end{equation}
La matrice $A$ permette quindi di mettere in relazione le matrici $\gamma^\mu$ con le loro aggiunte hermitiane. Si definiscono inoltre due matrici, $B$ e $C$, che permettono rispettivamente di mettere in relazione le matrici $\gamma^\mu$ con le loro coniugate complesse e con le loro trasposte, per cui
\begin{equation}
    (\gamma^\mu)^* = -t_0t_1B\gamma^\mu B^{-1},
\end{equation}
\begin{equation}
    (\gamma^\mu)^T = +t_0t_1C\gamma^\mu C^{-1},
\end{equation}
dove $t_0,t_1 = \pm 1$ e le matrici soddisfano le relazioni
\begin{equation}
    BB^\dag = CC^\dag = \bb{I},\quad C^T = -t_0C,\quad C^* = -t_0C^{-1}.
\end{equation}
Le matrici $B$ e $C$ appartengono alla classe degli \emph{intertwining operators} poiché mettono in relazione rappresentazioni associate alle matrici $\gamma$ e alle loro trasformate.

Nell'algebra di Clifford delle matrici $2^m \times 2^m$, sia per $D = 2m$ che per $D = 2m + 1$, si può distinguere tra matrici simmetriche ed anti-simmetriche, dove la proprietà di simmetria è definita nel modo seguente. Si consideri la famiglia di matrici $\{\Gamma^A\}$, applicando la matrice $C$ ad una matrice $\Gamma^A \in \{\Gamma^A\}$ si ottiene una matrice $C\Gamma^A$ che può essere simmetrica od anti-simmetrica, ma ciò dipende unicamente dal rango $r$ di $\Gamma^A$. Pertanto, si può scrivere 
\begin{equation}
\label{eq: condizione simmetrica matrici Cgamma}
    (C\Gamma^{(r)})^T = -t_r(C\Gamma^{(r)}),\quad t_r = \pm 1,
\end{equation}
dove $\Gamma^{(r)}$ e una matrice di rango $r$ in $\{\Gamma^A\}$. In particolare, se $t_r = +1$, la matrice $C\Gamma^{(r)}$ è simmetrica; se invece $t_r = -1$, è antisimmetrica. Per $r > 2$, i segni $t_r$ sono determinati da $t_0$ e $t_1$. Per i ranghi $r = 0$ e $r = 1$, da \eqref{eq: condizione simmetrica matrici Cgamma} si ottiene
\begin{equation}
    C^T = -t_0C,\quad (\gamma^\mu)^T = t_0t_1C\gamma^\mu C^{-1},
\end{equation}
si possono determinare le simmetrie di tutte le matrici $C\gamma^{\mu_1\cdots\mu_r}$ e quindi tutti i coefficienti $t_r$. Ad esempio, si trova che $r_2 = -t_0$, $t_3 = -t_1$ e che $t_{r + 4} = t_r$. In generale, i valori di $t_0$ e $t_1$, e quindi di tutti i $t_r$, dipendono dalla dimensione dello spaziotempo $D\mod{8}$ e dal rango $r\mod{4}$
\begin{ex}
    Per trovare $t_2$ si calcola
    \begin{equation}
    \begin{split}
        (C\gamma^{\mu\nu})^T &= \frac{1}{2}(C\gamma^\mu\gamma^\nu - C\gamma^\nu\gamma^\mu) \\
        &= \frac{1}{2}[(\gamma^\nu)^T(\gamma^\mu)^T - (\gamma^\mu)^T(\gamma^\nu)^T]C^T \\
        &= -\frac{t_0}{2}C(\gamma^\nu\gamma^\mu - \gamma^\mu\gamma^\nu) \\
        &= t_0C\gamma^{\mu\nu}.
    \end{split}
\end{equation}
Pertanto, da \eqref{eq: condizione simmetrica matrici Cgamma} si vede che $t_2 = -t_0$.
\end{ex}
Le matrici $C$ possono essere costruite esplicitamente mediante prodotti tensoriali di matrici di Pauli. Una possibile scelta è
\begin{equation}
\label{eq: definizione matrice C+}
    C_+=\sigma_1\otimes\sigma_2\otimes\sigma_1
    \otimes\sigma_2\otimes\cdots,
    \quad t_0t_1=+1,
\end{equation}
oppure
\begin{equation}
\label{eq: definizione matrice C-}
    C_-=\sigma_2\otimes\sigma_1\otimes\sigma_2
    \otimes\sigma_1\otimes\cdots,
    \quad t_0t_1=-1.
\end{equation}
La scelta dipende dalla dimensione dello spaziotempo $D$ e dalla rappresentazione delle matrici $\gamma$ adottata. Per dimensioni $D$ pari entrambe le scelte sono possibili e si può passare da una scelta all'altra rinominando $C$ con $C\gamma_*$, a meno di un fattore di fase di normalizzazione. Per dimensioni $D$ dispari, la scelta di $C$ è unica, di nuovo, a meno di un fattore di fase. Dato che si fa uso di rappresentazioni hermitiane, che soddisfano la condizione \eqref{eq: hermiticità matrici di Dirac}, la proprietà di simmetria di una matrice $\gamma$ determina anche quella della sua complessa coniugata. Per vederlo, si definisce $B$ in funzione di $C$ e $A$ come
\begin{equation}
\label{eq: definizione matrice B}
    B = it_0C\gamma^0 = t_0CA.
\end{equation}
Si può dimostrare poi che 
\begin{equation}
\label{eq: complesso coniugato gamma^*}
    (\gamma^\mu)^* = -t_0t_1B\gamma^\mu B^{-1},
\end{equation}
\begin{equation}
    (\gamma^\mu)^\dag = -A\gamma^\mu A^{-1}.
\end{equation}
Inoltre, utilizzando la definizione di $B$ e le proprietà di $C$, si ottiene
\begin{equation}
    \begin{split}
    B^*B &=-it_0C^*\gamma^0\,it_0C\gamma^0 \\
    &= -t_0C^{-1}(\gamma^0)^2t_0C\gamma^0 \\
    &= -t_0C^{-1}B\gamma^0B^{-1}C\gamma^0(-t_0t_1) \\
    &= t_1(C^{-1}B)\gamma^0(C^{-1}B)^{-1}\gamma^0 \\
    &= t_1\bb{I}.
    \end{split}
\end{equation}
In $ D = 4$, usando la rappresentazione di Weyl, si può scegliere $B = \gamma^0\gamma^1\gamma^3$, che risulta reale e simmetrica, implicando $C = i\gamma^3\gamma^1$. Per $t_0 = 1$ e $t_1 = -1$, si ottiene che $B = B^*$ e che $BB^* \equiv B^2 = \IM$. I segni associati ai prodotti anti-simmmetrizzati delle matrici $\gamma$ risultano quindi
\begin{equation}
    t_0 = +1,\quad t_1 = -1,\quad t_2 = -1,\quad t_3 = 1,\quad \dots\quad.
\end{equation}
Con una scelta alternativa delle matrici di $B$ e $C$ si ottiene invece
\begin{equation}
    t_0 = +1,\quad t_1 = +1,\quad t_2 = -1,\quad t_3 = -1,\quad \dots\quad.
\end{equation}
Questi risultati mostrano che le proprietà di simmetria dei prodotti di matrici gamma dipendono dalla scelta degli operatori intreccianti e dalla dimensione dello spaziotempo.

Le proprietà \eqref{eq: condizione simmetrica matrici Cgamma} e \eqref{eq: complesso coniugato gamma^*} valgono per la rappresentazione \eqref{eq: costruzione matrici di Dirac euclidee} usando le matrici \eqref{eq: definizione matrice C+}, \eqref{eq: definizione matrice C-} e \eqref{eq: definizione matrice B}. In un'altra rappresentazione, legata da \eqref{eq: matrici gamma legate da coniugazione}, le matrici $C$ e $B$ sono date da 
\begin{equation}
    C' = (S^{-1})^TCS^{-1},\quad B' = (S^{-1})^TBS^{-1}.
\end{equation}

% \subsubsection{Simmetrie delle matrici \texorpdfstring{$\gamma$}{gamma}}
% Se $\b{\Psi} = \Psi^\dag\gamma^0$, si definisce $A = i\gamma^0$ tale che 
% \begin{equation}
%     A\gamma^\mu A^{-1} = -(\gamma^\mu)^\dag \quad\Longleftrightarrow\quad A\gamma^\mu = -(\gamma^\mu)^\dag A,
% \end{equation}
% con il quale quindi si passa da $\gamma$ ad $\gamma^\dag$. Si definisce poi $B$ come la mappa che permette di passare da $\gamma$ a $\gamma^*$ e con $C$ quella tra $\gamma$ e $\gamma^T$, come 
% \begin{equation}
%     (\gamma^\mu)^* = -t_0t_1B\gamma^\mu B^{-1},
% \end{equation}
% \begin{equation}
%     (\gamma^\mu)^T = +t_0t_1C\gamma^\mu C^{-1},
% \end{equation}
% dove $t_0 = t_1 = \pm 1$ e con $BB^\dag = CC^\dag = \IM$ e con $C^T = -t_0C$ o $C^* = -t_0C^{-1}$. Si definiscono quindi le $\Gamma^{(\nu)} \in \{\IM,\gamma^\mu,\gamma^{\mu\nu},\dots\}$, tale che 
% \begin{equation}
%     (C\Gamma^{(\nu)})^T = t_r{(C\Gamma^{(r)})},\quad t_r = \pm 1.
% \end{equation}
% Quindi, $t_r$ con $r > 2$, è fissato da $t_0$ e $t_1$. Ad esempio
% \begin{equation}
%     \begin{split}
%         (C\gamma^{\mu\nu})^T &= \frac{1}{2}(C\gamma^\mu\gamma^\nu - C\gamma^\nu\gamma^\mu) \\
%         &= \frac{1}{2}[(\gamma^\nu)^T(\gamma^\mu)^T - (\gamma^\mu)^T(\gamma^\nu)^T]C^T \\
%         &= -\frac{t_0}{2}C(\gamma^\nu\gamma^\mu - \gamma^\mu\gamma^\nu) \\
%         &= t_0C\gamma^{\mu\nu}.
%     \end{split}
% \end{equation}
% Da qui si trova che $t_2 = -t_0$, si può poi trovare che $t_3 = -t_1$, $t_{r + 4} = t_r$ e così via. Si può poi costruire esplicitamente una rappresentazione di $C$ come
% \begin{equation}
%     C_+ = \sigma_1 \ot \sigma_2 \ot \sigma_1 \ot \sigma_2 \ot \cdots,\quad t_0t_1 = 1,
% \end{equation}
% o
% \begin{equation}
%     C_- = \sigma_2 \ot \sigma_1 \ot \sigma_2 \ot \sigma_1 \ot \cdots,\quad t_0t_1 = -1,
% \end{equation}
% a differenza di $A = i\gamma^0$ che non dipende dalla segnatura. In modo simile, si ha che 
% \begin{equation}
%     B = it_0C\gamma^0 = t_0CA.
% \end{equation}
% Si dimostra 
% \begin{equation}
%     (\gamma^\mu)^\dag = (\gamma^T)^* = \cdots = -\gamma^0\gamma^\mu\gamma^0 = -A\gamma^\mu A^{-1}.
% \end{equation}
% Si dimostra
% \begin{equation}
%     \begin{split}
%         B^*B &= -it_0C^*\gamma^*it_0C\gamma^0 \\
%         &= -t_0C^{-1}(\gamma^0)^*C\gamma^0 \\
%         &= -t_0C^{-1}B\gamma^0B^{-1}C\gamma^0(-t_0t_1) \\
%         &= t_1(C^{-1}B)\gamma^0(C^{-1}B)^{-1}\gamma^0 \\
%         &= -t_1.
%     \end{split}
% \end{equation}
% In $D = 4$, da $B = \gamma^0\gamma^1\gamma^3$ e $C = i\gamma^3\gamma^1$ so ha $t_0 = 1$, allora $B = B^*$ e $BB^* = B^2 = \cdots = \IM$ e $t_1 = -1$, scegliendo quindi 
% \begin{equation}
%     t_0 = 1,t_1 = -1,t_2 = -1,t_3 = 1,\dots,
% \end{equation}
% invece che
% \begin{equation}
%     t_0 = 1,t_1 = 1,t_2 = -1,t_3 = -1,\dots.
% \end{equation}

\section{Spinori in dimensione generica}
In equazione \eqref{eq: aggiunto di Dirac} si è definito l'aggiunto di Dirac, nella quale definizione compariva il complesso coniugato dello spinore $\Psi(x)$, e lo si è usato per definire una forma bilineare invariante di Lorentz così da costruire un'azione per il campo spinoriale di Dirac. Si definisce adesso \emph{l'aggiunto di Majorana} per qualsiasi spinore $\lambda$ usando il suo trasposto e la matrice $C$ introdotta più sopra come 
\begin{equation}
\label{eq: aggiunto di Majorana}
    \b{\lambda} := \lambda^TC.
\end{equation}
La forma bilineare $\b{\lambda}\chi = \lambda^TC\chi$ è un invariante di Lorentz. La variazione di $\b{\lambda}\chi$ sotto una trasformazione di Lorentz infinitesima può essere scomposta in
\begin{equation}
\delta(\b{\lambda}\chi) = (\delta\b{\lambda})\chi + \b{\lambda}\delta\chi,
\end{equation}
dove le singole variazioni risultano essere
\begin{equation}
    \delta\lambda^T = -\frac{1}{4}\lambda^T(\gamma_{\mu\nu})^T\lambda^{\mu\nu},\quad \delta\chi = -\frac{1}{4}\lambda^{\mu\nu}\gamma_{\mu\nu}\chi,
\end{equation}
per cui si ha 
\begin{equation}
    \delta(\b{\lambda}\chi) = -\frac{1}{4}\lambda^{\mu\nu}[\lambda^T(\gamma_{\mu\nu})^TC\chi + \lambda^TC\gamma_{\mu\nu}\chi] = 0
\end{equation}
da cui si definiscono gli invarianti $\lambda^T(\gamma_{\mu\nu})^TC\chi$ e $\lambda^TC\gamma_{\mu\nu}\chi$.

\paragraph{Numeri di Grassmann.} I numeri di Grassmann sono definiti come quantità che anti-commutano con altri numeri di Grassmann, ma sono tali da commutare con i numeri ordinari. Dati due numeri di Grassmann $\theta$ e $\eta$ ed un numero ordinario $x$, si hanno
\begin{equation}
    \begin{cases}
        \theta\eta = -\eta\theta \\
        \theta x = x\theta
    \end{cases}.
\end{equation}
I numeri di Grassmann son nilpotenti, ossia
il quadrato di un numero di Grassmann $\theta$ è nullo, cioè $\theta\theta = 0$. Considerando ancora i due numeri di Grassmann $\theta$ ed $\eta$, questi soddisfano la relazione
\begin{equation}
    (\theta\eta)^* = \eta^*\theta^* = -\theta^*\eta^*.
\end{equation}
Data una funzione lineare $f(\theta)$ del tipo $f(\theta) = a + b\theta$, per distinguere i numeri di Grassmann dai numeri complessi ordinari, si chiamano questi ultimi \emph{c-numeri} dove la ``c'' sta per ``commutanti''. A meno che non sia specificato diversamente, si assume nel seguito che le componenti degli spinori siano numeri anti-commutanti di Grassmann. Questo riflette l'importante relazione fisica tra lo spin e la statistica.

Usando la definizione \eqref{eq: aggiunto di Majorana} e la proprietà \eqref{eq: condizione simmetrica matrici Cgamma}, si ottiene l'utile relazione
\begin{equation}
\label{eq: Majorana flip}
    \b{\lambda}\gamma_{\mu_1\cdots\nu_r}\chi = t_r\b{\chi}\gamma_{\mu_1\cdots\mu_r}\lambda.
\end{equation}
Il segno meno risultante dallo scambio dell'ordine dei degli spinori con numeri di Grassmann come componenti è stato incorporato. La proprietà di simmetria in \eqref{eq: Majorana flip} è valida per gli spinori di Dirac, ma l'applicazione principale in supergravità è legata ad un particolare tipo di spinori, cioè gli \emph{spinori di Majorana}. Per questo motivo, ci si riferisce alla relazione \eqref{eq: Majorana flip} con il nome di \emph{Majorana flip}. 

Spesso si incontra il caso in cui il bilineare contiene il prodotto di singole matrici $\gamma$. In questi casi, si può dimostrare che per le \emph{dimensioni di Majorana} $D = 2,3,4\mod{8}$ si ha 
\begin{equation}
    \b{\lambda}\gamma^{\mu_1}\cdots\gamma^{\mu_p}\chi = (-1)^p\b{\chi}\gamma^{\mu_p}\cdots\gamma^{\mu_1}\lambda.
\end{equation}
Le relazioni precedenti implicano anche la seguente regola. Per qualsiasi relazioni tra spinori che includono le matrici $\gamma$, esiste una relazione di corrispondenza tra gli spinori barrati,
\begin{equation}
    \chi_{\mu_1\cdots\mu_r} = \gamma_{\mu_1\cdots\mu_r}\lambda \quad\Longrightarrow\quad \b{\chi}_{\mu_1\cdots\mu_r} = t_0t_r\b{\lambda}\gamma_{\mu_1\cdots\mu_r}.
\end{equation}

In dimensioni pari si definiscono le parti left e right degli spinori usando i proiettori \eqref{eq: proiettori di chiralità}. La definizione \eqref{eq: aggiunto di Majorana} implica che la chiralità dello spinore coniugato dipende da $t_0t_D$, e si ottiene 
\begin{equation}
    \chi = P_L\lambda \quad \longrightarrow \quad \b{\chi} = 
    \begin{cases}
        \b{\lambda}P_L, &D = 0,4,8,\dots \\
        \b{\lambda}P_R, &D = 2,6,10,\dots 
    \end{cases}.
\end{equation}
% Questa si può verificare ricordando che $s^T = s$ con $S \in \bb{C}$, per cui 
% \begin{equation}
%     \begin{split}
%         (\b{\lambda}\gamma_{\mu_1\cdots\mu_r}\chi)^T &= \b{\lambda}\gamma_{\mu_1\cdots\mu_r}\chi \\
%         &= -\chi^T(\gamma_{\mu_1\cdots\mu_r})^TC^T\lambda,
%     \end{split}
% \end{equation}
% Se $\lambda = \chi$ e $t_r = -1$, allora la quantità $\b{\lambda}\gamma^{\mu\nu}\lambda$ si annulla in $D = 4$, ma in generale dipende dalla dimensione. 
Per scambiare di posto un numero di spinori maggiore di due sono necessarie delle accortezze. Per un prodotto di quattro spinori ad esempio, il caso più comune, si può usare una forma dell'identità di Fierz. Siano $\lambda_1$, $\lambda_2$, $\lambda_3$ e $\lambda_4$ quattro spinori generici, l'identità di Fierz permette di scrivere
\begin{equation}
\label{eq: identità di Fierz 4 spinori}
    (\b{\lambda}_1\lambda_2)(\b{\lambda}_3\lambda_4) = \nsum_A\frac{1}{2^m}\b{\lambda}_1\Gamma^A\lambda_4\b{\lambda}_3\Gamma^A\lambda_2,
\end{equation}
dove, in dimensione pari,
\begin{equation}
    \Gamma^A \in \{\Gamma^A\} = \{\IM,\gamma^\mu,\dots,\gamma^{\mu_1\cdots\mu_D}\} \in \alg{Cliff}{2m}
\end{equation}
oppure, in dimensione dispari,
\begin{equation}
    \Gamma^A \in \{\Gamma^A\} = \{\IM,\gamma^\mu,\dots,\gamma^{\mu_1\cdots\mu_m}\} \in \alg{Cliff}{2m + 1}.
\end{equation}
La relazione \eqref{eq: identità di Fierz 4 spinori} può poi essere generalizzata per includere matrici generali $M$, $M'$ dell'algebra di Clifford. Quando gli spinori $\lambda_1$, $\lambda_2$, $\lambda_3$ e $\lambda_4$ non sono del tutto indipendenti, può spesso capire che alcuni termini della somma si annullino a causa di relazioni di simmetria come \eqref{eq: Majorana flip}. 

\subsection{Coniugazione di carica}
Fino ad ora non si è ancora discusso della complessa coniugazione dei campi spinoriali che si sta studiando. La complessa coniugazione è necessaria per verificare, ad esempio, che un termine nella lagrangiana comprendente dei bilineari di spinori sia hermitiano o meno. Tuttavia, la complessa coniugazione di un bilineare è un'operazione delicata dato che sono coinvolte le hermiticità della matrice $C$ in \eqref{eq: aggiunto di Majorana} e delle matrici $\gamma$. Pertanto, si introduce l'operazione di \emph{coniugazione di carica} più semplice da usare nei calcoli. Per ogni scalare, definito come una quantità i cui indici spinoriali sono totalmente contratti, la coniugazione di carica e la complessa coniugazione sono equivalenti. Dato che la lagrangiana è uno scalare, la coniugazione di carica può essere usata per manipolare i termini che la formano. 

Si definisce prima di tutto il \emph{coniugato di carica} di un generico spinore $\lambda$ come 
\begin{equation}
\label{eq: coniugato di carica}
    \lambda^C := B^{-1}\lambda^*.
\end{equation}
Il coniugato di carica barrato di un generico spinore $\lambda$ è pertanto definito, usando \eqref{eq: aggiunto di Majorana} e \eqref{eq: definizione matrice B}, come\footnote{Facendo il coniugato di Majorana di \eqref{eq: coniugato di carica} si ha $\bigb{\lambda^C} = (B^{-1}\lambda^*)^TC = \lambda^\dag (B^{-1})^TC$, da cui, usando che $(B^{-1})^T = -t_1B^{-1}$, si trova \eqref{eq: aggiunto del coniugato di carica}.} 
\begin{equation}
\label{eq: aggiunto del coniugato di carica}
    \bigb{\lambda^C} = (-t_0t_1)i\lambda^\dag\gamma^0.
\end{equation}
\begin{obs}
    Si noti che \eqref{eq: aggiunto del coniugato di carica} corrisponde all'aggiunto di Dirac definito in \eqref{eq: aggiunto di Dirac} a meno del fattore numerico $-t_0t_1$. Il significato di ciò sarà più chiaro quando si introdurranno gli spinori di Majorana. Si noti anche che $-t_0t_1 = 1$ in $D = \{2,3,4,10,11\}$.
\end{obs}
Il coniugato di carica di qualsiasi matrice $M \in \alg{Cliff}{D}$ di dimensione $2^m \times 2^m$ è definito come
\begin{equation}
\label{eq: coniugato di carica matrice M di clifford generica}
    M^C := B^{-1}M^*B.
\end{equation}
La coniugazione di carica non cambia l'ordine delle matrici, cioè $(MN)^C = M^CN^C$. Il tipo di matrici $M$ con cui si ha maggiormente a che fare sono le matrici $\gamma^\mu$, per cui sono particolarmente importanti le proprietà di coniugazione delle matrici generatrici $\gamma$, che sono 
\begin{equation}
    \gamma_\mu^C = B^{-1}\gamma_\mu^*B = -t_0t_1\gamma_\mu.
\end{equation}
In particolare, per la matrice di chiralità si ha
\begin{equation}
    \gamma_*^C = (-1)^{(D/2) + 1}\gamma_*.
\end{equation}
Con queste relazioni, si può derivare la regola per la complessa coniugazione di un bilineare spinoriale che include una matrice arbitraria $M$. Questa è 
\begin{equation}
\label{eq: coniugato di carica di un bilineare}
    (\b{\chi}M\lambda)^* \equiv (\chi^TCM\lambda)^C = -t_0t_1\bigb{\chi^C}M^C\lambda^C.
\end{equation}
Una scambio nell'ordine dei campi fermionici è necessaria per la derivazione\footnote{+DERIVAZIONE}, ma non c'è nessuna variazione nell'ordine del risultato finale per il coniugato di carica di qualsiasi bilineare. Si può quindi considerare la relazione \eqref{eq: coniugato di carica di un bilineare} come la proprietà di coniugazione di carica appropriata quando il coniugato di uno spinore è definito come in \eqref{eq: aggiunto di Majorana}. 
\begin{obs}
    Si può dimostrare che per ogni spinore $(\lambda^C)^C = -t_1\lambda$, mentre per ogni matrice $(M^C)^C = M$.
\end{obs}

% Si calcola quindi
% \begin{equation}
%     \begin{split}
%         (\b{\chi}M\lambda)^* \equiv (\chi^TCM\lambda)^C &= (\chi_a(CM)_{ab}\lambda_b)^* \\
%         &= \lambda_b^*(CM)^*_{ab}\chi_a^* \\
%         &= \chi_a^*(CM)^*_{ab}\lambda_b^* \\
%         &= -\chi^TC^*M^*\lambda^* \\
%         &= i\chi^\dag\gamma^0 M^C \lambda^C,
%     \end{split}
% \end{equation}
% usando poi la definizione di $M^C$ in \eqref{eq: coniugato di carica matrice M di clifford generica} insieme alle relazioni $C^* = -t_0C^{-1}$ e $B = it_0C\gamma^0$, si ottiene la regola
% \begin{equation}
%     (\b{\chi}M\lambda)^* \equiv (\chi^TCM\lambda)^C = -t_0t_1\bigb{\chi^C}M^C\lambda^C.
% \end{equation}
% Dunque, l'operazione di coniugazione di carica agisce come il complesso coniugato a meno di un segno. 

\section{Spinori di Majorana}
Gli spinori di Dirac in generale hanno componenti complesse. Si vuole capire se è possibile imporre delle condizioni affinché le entrate siano invece reali. Una prima condizione che si potrebbe pensare di imporre è
\begin{equation}
    f - f^* = 0,
\end{equation}
dove $f$ è una qualsiasi funzione, così da avere $\Psi = \Psi^*$. Tuttavia, così facendo lo spinore non soddisfa le trasformazioni di Lorentz. Uno spinore che invece soddisfa la condizione di realtà delle componenti e le trasformazioni di Lorentz è il coniugato di carica
\begin{equation}
    \Psi^c = B^{-1}\Psi^*,
\end{equation}
con il quale si può imporre che $\Psi = \Psi^c$, equivalente a chiedere $\Psi^* = B\Psi$. Infatti, da $\Psi^* = B\Psi$ si ha 
\begin{equation}
    \Psi = B^*\Psi^* = B^*B\Psi,
\end{equation}
da cui si deve richiedere che $B^*B = \IM$, ma allora $t_1 = -1$. Possibile solo in alcune dimensioni, $D = 2,3,4 \mod{8}$, dove si ha che il coniugato di Majorana corrisponde all'aggiunto di Dirac, per cui 
\begin{equation}
    \b{\Psi^c} = \Psi^\dag A.
\end{equation}
\begin{obs}
    Ci sono rappresentazioni dove le matrici $C$ sono completamente reali (\emph{really-real representations}). Quando non esistono gli spinori di Majorana, si può introdurre una condizione pseudo-majorana. 
\end{obs}
Solo in dimensione $D = 2\mod{8}$ esistono gli spinori di Majorana che sono anche spinori di Weyl, detti \emph{spinori di Majorana-Weyl}. Per $D = 4$, non si può avere oggetti di questo tipo, ma si può avere o uno o l'altro caso, sia quindi $\Psi$ uno spinore di Majorana, in tal caso
\begin{equation}
    \Psi = (P_L + P_R)\Psi = P_L\Psi + P_R\Psi,
\end{equation}
dove $P_L\Psi$ e $P_R\Psi$ sono singolarmente spinori di Weyl, anche se la loro somma è uno spinore di Majorana. Si può vedere in particolare che 
\begin{equation}
    (P_L\Psi)^c = P_L^c\Psi^c = P_L^c\Psi = \frac{1}{2}(\IM + \gamma_*^c)\Psi,
\end{equation}
dove 
\begin{equation}
    \gamma_*^c = B^{-1}\gamma_*^*B = -\gamma_*,
\end{equation}
per cui
\begin{equation}
    P_L^c = \frac{1}{2}(\IM - \gamma_*) \equiv P_R,
\end{equation}
dunque 
\begin{equation}
    (P_L\Psi)^c = P_R\Psi,
\end{equation}
e ovviamente
\begin{equation}
    P_L\Psi = (P_R\Psi)^c.
\end{equation}
Quindi, la coniugazione di carica permette di passare da uno spinore all'altro, confermando la natura di spinore di Weyl dei singoli spinori. 

L'azione di Dirac era stata definita come 
\begin{equation}
    S_{Dirac} = -\int d^4x\,\b{\Psi}(\sl{\p} - m)\Psi,
\end{equation}
ma se $\Psi$ è uno spinore di Majorana la forma dell'azione non cambia dato che coniugato di Majorana corrisponde all'aggiunto di Dirac; quindi semplicemente
\begin{equation}
    S_{Majorana} = -\frac{1}{2}\int d^4x\,\b{\Psi}(\sl{\p} - m)\Psi.
\end{equation}
Tuttavia, si deve controllare se gli invarianti in $S_{Majorana}$ sono effettivamente tali. Se si suppone che i campi siano commutati, il Majorana flip restituisce che 
\begin{equation}
    \b{\Psi}\Psi = -\b{\Psi}\Psi,
\end{equation}
rendendo quindi il termine di massa in $S_{Majorana}$ nullo. Anche se questo di per se non sarebbe un problema dato che basterebbe richiedere che $m = 0$, il fatto è che anche il termine cinetico si annulla se i campi sono commutanti; infatti
\begin{equation}
    \b{\Psi}\gamma^\mu\p_\mu\Psi = \p_\mu\b{\Psi}\gamma^\mu\Psi\quad\Longrightarrow\quad \b{\Psi}\gamma^\mu\p_\mu\Psi = \frac{1}{2}\p_\mu(\b{\Psi}\gamma\Psi),
\end{equation}
che essendo un termine di bordo non contribuisce all'azione. In sintesi, se i campi sono commutanti l'azione è nulla, invece se questi anti-commutano, allora l'azione non si annulla e sono presenti sia il termine cinetico che quello di massa. Calcolando la variazione $\delta S$, si ha 
\begin{equation}
    \b{\Psi}\delta\Psi = \delta\b{\Psi}\Psi
\end{equation}
e una relazione simile con le $\gamma^\mu$ incluse, per cui 
\begin{equation}
    \delta S = -\int d^4x\,\delta\b{\Psi}(\gamma^\mu\p_\mu - m)\Psi,
\end{equation}
da cui si ricava che lo spinore di Majorana soddisfa l'equazione di Dirac
\begin{equation}
    (\sl{\p} - m)\Psi = 0.
\end{equation}
Inoltre, si può scrivere l'azione di Majorana in funzione degli spinori di Weyl, 
\begin{equation}
    \begin{split}
        S_{Majorana} &= -\frac{1}{2}\int d^4x\,\b{\Psi}(\sl{\p} - m)(P_L + P_R)\Psi \\
        &= -\frac{1}{2}\int d^4x\,\Psi^TC(P_L + P_R)(\sl{\p} - m)(P_L + P_R)\Psi \\
        &= -\frac{1}{2}\int d^4x\,[2\b{\Psi}\Psi\gamma^\mu\p_\mu P_L\Psi - m\b{\Psi}P_L\Psi - m\b{\Psi}P_R\Psi] \\
        &\equiv S[P_L\Psi,(P_L\Psi)^c].
    \end{split}
\end{equation}

\subsection{Simmetrie}
Sia $\Psi$ uno spinore di Dirac. Non è molto difficile dimostrare che se si applica la trasformazione 
\begin{equation}
    \Psi \to \Psi' \equiv e^{i\theta}\Psi,\quad \theta \in \bb{R},
\end{equation}
l'azione è invariante. Questa trasformazione è un esempio di simmetria interna globale, in particolare una simmetria del gruppo $U(1)$. Se invece lo spinore $\Psi$ è uno spinore di Majorana, la trasformazione precedente non vede l'azione rimanere invariante, si può però operare la trasformazione
\begin{equation}
    \Psi \to \Psi' \equiv e^{i\theta\gamma_*}\Psi,
\end{equation}
che conserva la condizione di Majorana,
\begin{equation}
    B^{-1}\Psi^* \to B^{-1}e^{-i\theta\gamma_*^*}\Psi^* = B^{-1}e^{-i\theta\gamma_*^*}BB^{-1}\Psi^*,
\end{equation}
se $BB^{-1} = \IM$.
% DA VEDERE
% da cui, se $BB^{-1} = \IM$, si ha
% \begin{equation}
%     B^{-1}e^{-i\theta\gamma_*^*}B = \IM.
% \end{equation}
Si vede che 
\begin{equation}
    \begin{split}
        B^{-1}e^{-i\theta\gamma_*^*} &= B^{-1}[\IM - i\theta\gamma_*^* + \frac{1}{2}(-i)^2\theta^2(\gamma_*^*)^2 + \cdots] \\
        &= e^{i\theta\gamma_*},
    \end{split}
\end{equation}
dove si è usato che $(\gamma_*^*)^2 = \gamma_*^*BB^{-1}\gamma_*^*$ e per ogni scambio di $\theta$ con $\gamma_*$ c'è un segno. Si può far vedere che lo spinore $\b{\Psi}$ invece trasforma come 
\begin{equation}
    \b{\Psi} \equiv \Psi^TC \to \Psi^Te^{i\theta\gamma_*^T}C = \Psi^TCC^{-1}e^{i\theta\gamma_*^T}C = \b{\Psi}e^{i\theta\gamma_*}.
\end{equation}
Per quanto riguarda $\b{\Psi}\Psi$ invece, si può far vedere che non è invariante sotto tale trasformazione dato che gli esponenziali non si eliminano. Ponendo $m = 0$, il fattore cinetico diventa invece
\begin{equation}
    \b{\Psi}\gamma^\mu\p_\mu\Psi \to \b{\Psi}e^{i\theta\gamma_*}\gamma^\mu e^{i\theta\gamma_*}\p_\mu\Psi,
\end{equation}
ma se $\acomm{\gamma_*}{\gamma^\mu} = 0$, allora espandendo gli esponenziali e scambiando la $\gamma_*$, si ha 
\begin{equation}
    \b{\Psi}\gamma^\mu\p_\mu\Psi \to \b{\Psi}e^{i\theta\gamma_*}e^{-i\theta\gamma_*}\sl{\p}\Psi = \b{\Psi}\sl{\p}\Psi,
\end{equation}
che risulta quindi essere invariante. Da qui si giustifica la necessità di porre $m = 0$ nell'azione di Majorana. La corrente conservata è
\begin{equation}
    J^\mu_* = \frac{1}{2}\b{\Psi}\gamma^\mu\gamma_*\Psi,
\end{equation}
la quale soddisfa 
\begin{equation}
    \p_\mu J^\mu_* = -im\b{\Psi}\gamma_*\Psi,
\end{equation}
che è nulla appunto se $m = 0$.